% ============================================================
% Week 03 section --- NLP for Structured HAZOP Reasoning
%
% Research planning and literature synthesis:
%   - Single et al. (2020)
%   - Mao et al. (2020)
%   - HAZOPCT / Li & Zhao (2025)
%   - Lee et al. (2026)
%
% This week moves from paper-by-paper analysis toward defining
% the case study, research questions, causal representation,
% experimental design, and validation strategy.
% ============================================================

\chapter[Literature Synthesis and Research Design]
{Literature Synthesis and Research Design}

\section{Objective}

The objective of this week was to consolidate the literature studied during
the previous weeks and use it to begin defining the research problem, case
study, and experimental methodology for the project. Rather than treating
the individual papers as isolated approaches, the analysis focuses on the
different stages involved in automated HAZOP analysis: knowledge
representation, information extraction, knowledge-graph construction,
causal reasoning, missing-relationship inference, and direct HAZOP
generation.

The purpose of this synthesis is not yet to claim a finalized research gap.
Instead, it is to determine which parts of the HAZOP reasoning pipeline have
already been addressed, what assumptions existing approaches make, and which
questions should be investigated through a focused case study.

\begin{glance}[Week 3 at a glance]
This week moves from reviewing individual automation approaches to defining
the project's research design. The central working direction is to extract
causal propositions from historical HAZOP reports, represent them explicitly,
and evaluate whether the resulting structure supports multi-hop reasoning.
The case study, schema, and ground-truth strategy remain provisional until
candidate datasets have been examined.
\end{glance}

\section{From Automated HAZOP Generation to Structured Causal Reasoning}

The literature studied so far suggests that ``automating HAZOP'' is not a
single technical problem. Different approaches automate different stages of
the overall process.

Single et al.\ (2020) demonstrate that HAZOP knowledge can be represented
using an ontology and that symbolic inference can be used to generate
multi-step HAZOP scenarios. Their system operates on an already structured
process and safety representation, with expert knowledge encoded in the
ontology before the reasoning process begins
\cite{single2020ontologyhazop}.

Mao et al.\ (2020) move the knowledge-acquisition problem closer to the
document level by developing a process-safety knowledge graph from
heterogeneous process-safety information. Their approach uses natural
language processing for named-entity recognition and relation extraction
from textual sources, including HAZOP reports. The resulting graph is then
used for applications including information retrieval and causal-chain
extraction \cite{mao2020processsafetykg}.

Li and Zhao's HAZOPCT (2025) extends the knowledge-graph direction toward
HAZOP completeness. Their approach constructs a HAZOP knowledge graph from
multiple sources and applies graph-based learning to predict potentially
missing relationships, allowing existing HAZOP information to be used for
the identification of additional conclusions
\cite{lizhao2025hazopct}.

Lee et al.\ (2026) investigate a different direction: direct generation of
complete HAZOP worksheets using multimodal large language models. Their
experiments show that LLM-generated worksheets can achieve high textual
similarity to expert-prepared worksheets, while the proportion of semantically
valid generated scenarios remains considerably lower. This distinction is
important because linguistic similarity alone does not establish the
correctness of a causal or process-safety argument
\cite{lee2026llmhazop}.

The four approaches therefore address different parts of the overall problem,
as summarized in Figure~\ref{fig:w3-literature-landscape}.

\begin{figure}[htbp]
\centering
\begin{tikzpicture}[
    node distance=6mm and 7mm,
    box/.style={
        draw=hzblue,
        fill=hzbluebg,
        rounded corners,
        align=center,
        text width=31mm,
        minimum height=13mm,
        font=\scriptsize
    },
    arrow/.style={ar}
]

    \node[box] (single)
        {Single et al.\\(2020)\\
         Ontology + reasoning};

    \node[box, right=of single] (mao)
        {Mao et al.\\(2020)\\
         NLP + knowledge graph};

    \node[box, below=of single] (hazopct)
        {HAZOPCT\\(2025)\\
         KG completion};

    \node[box, right=of hazopct] (lee)
        {Lee et al.\\(2026)\\
         LLM generation};

    \node[below=8mm of $(hazopct)!0.5!(lee)$,
          align=center,
          text width=72mm,
          font=\small\itshape]
        (synthesis)
        {Different stages of automated HAZOP analysis};

    \draw[arrow] (single) -- (mao);
    \draw[arrow] (hazopct) -- (lee);

\end{tikzpicture}
\caption[Literature landscape for automated HAZOP analysis]
{The papers studied so far address different stages of automated HAZOP
analysis: structured symbolic reasoning, knowledge extraction and graph
construction, graph-based completion, and direct LLM-based generation.}
\label{fig:w3-literature-landscape}
\end{figure}

\section{A Layered View of the HAZOP Problem}

The comparison indicates that several capabilities must be distinguished when
discussing automated HAZOP analysis.

\begin{enumerate}
    \item \textbf{Information extraction:} identifying relevant entities,
    events, and relationships from documents.

    \item \textbf{Representation:} encoding extracted information using a
    structured representation such as a table, ontology, knowledge graph, or
    causal graph.

    \item \textbf{Causal reasoning:} determining how deviations, causes,
    effects, consequences, and safeguards are related and how events
    propagate through intermediate states.

    \item \textbf{Inference and validation:} identifying potentially missing
    relationships or scenarios and determining whether the resulting
    reasoning is consistent with process and safety knowledge.
\end{enumerate}

These capabilities are related but are not interchangeable. In particular,
identifying the concepts ``cooling failure'', ``high temperature'', and
``pressure increase'' does not by itself establish that they form an ordered
causal sequence.

The distinction is illustrated in Figure~\ref{fig:w3-layered-problem}.

\begin{figure}[htbp]
\centering
\begin{tikzpicture}[
    node distance=6mm,
    box/.style={
        draw=hzblue,
        fill=hzbluebg,
        rounded corners,
        align=center,
        text width=48mm,
        minimum height=10mm,
        font=\scriptsize
    },
    arrow/.style={ar}
]

    \node[box] (text)
        {HAZOP report text};

    \node[box, below=of text] (extract)
        {Information extraction\\
         entities and relations};

    \node[box, below=of extract] (repr)
        {Structured representation\\
         events and typed relations};

    \node[box, below=of repr] (reason)
        {Causal reasoning\\
         multi-hop propagation};

    \node[box, below=of reason] (valid)
        {Inference and validation\\
         process-safety consistency};

    \draw[arrow] (text) -- (extract);
    \draw[arrow] (extract) -- (repr);
    \draw[arrow] (repr) -- (reason);
    \draw[arrow] (reason) -- (valid);

\end{tikzpicture}
\caption[Layered view of automated HAZOP reasoning]
{A layered view of the HAZOP problem. Information extraction is an
intermediate step; causal reasoning additionally requires relationships,
direction, ordering, and propagation between events.}
\label{fig:w3-layered-problem}
\end{figure}

This distinction is particularly relevant to the present project. A system
may correctly identify individual HAZOP concepts while still failing to
recover the reasoning connecting those concepts.

\section{Why Causal Structure Is the Central Focus}

The literature studied so far suggests that the causal structure of HAZOP
information is more important than the presence of individual keywords alone.

Consider the following hypothetical description:

\begin{center}
\emph{Cooling-water failure may increase reactor temperature, which can cause
pressure buildup and ultimately result in loss of containment.}
\end{center}

A conventional information-extraction system may identify:

\begin{center}
\textit{cooling-water failure}, \quad
\textit{reactor temperature}, \quad
\textit{pressure buildup}, \quad
\textit{loss of containment}.
\end{center}

However, identifying these concepts does not establish the causal structure:

\begin{center}
\[
\text{Cooling failure}
\rightarrow
\text{Temperature increase}
\rightarrow
\text{Pressure increase}
\rightarrow
\text{Loss of containment}.
\]
\end{center}

The direction and ordering of these relationships are essential. A system
that detects all four concepts but assigns incorrect relationships has not
recovered the causal reasoning expressed by the original statement.

This distinction is also consistent with the progression of the literature.
Single et al.\ demonstrate that HAZOP reasoning can contain multi-step and
branching causal structures, while Mao et al.\ demonstrate how process-safety
information can be organized into a knowledge graph
\cite{single2020ontologyhazop,mao2020processsafetykg}.

The present investigation therefore focuses on the connection between
natural-language information and an explicit causal representation.

\section{Current Working Research Pipeline}

Based on the literature studied so far, the current working hypothesis is to
investigate a pipeline that begins with historical HAZOP reports rather than
with an already constructed process ontology.

The proposed high-level pipeline is:

\begin{center}
\emph{Historical HAZOP reports}
$\rightarrow$
\emph{NLP/LLM extraction}
$\rightarrow$
\emph{causal information}
$\rightarrow$
\emph{structured causal representation}
$\rightarrow$
\emph{causal reasoning}
$\rightarrow$
\emph{validation}.
\end{center}

Figure~\ref{fig:w3-working-pipeline} presents the same idea in a compact
system representation.

\begin{figure}[htbp]
\centering
\begin{tikzpicture}[
    node distance=5mm and 4mm,
    box/.style={
        draw=hzblue,
        fill=hzbluebg,
        rounded corners,
        align=center,
        text width=25mm,
        minimum height=12mm,
        font=\scriptsize
    },
    arrow/.style={ar}
]

    \node[box] (reports)
        {Historical\\HAZOP reports};

    \node[box, right=of reports] (extract)
        {NLP / LLM\\
         causal extraction};

    \node[box, right=of extract] (graph)
        {Structured\\
         causal representation};

    \node[box, right=of graph] (reason)
        {Multi-hop\\
         reasoning};

    \node[box, right=of reason] (valid)
        {Process + expert\\
         validation};

    \draw[arrow] (reports) -- (extract);
    \draw[arrow] (extract) -- (graph);
    \draw[arrow] (graph) -- (reason);
    \draw[arrow] (reason) -- (valid);

\end{tikzpicture}
\caption[Current working research pipeline]
{Current working pipeline for investigating structured HAZOP reasoning from
historical reports. This is a working research hypothesis rather than a
finalized system architecture.}
\label{fig:w3-working-pipeline}
\end{figure}

The purpose of this pipeline is not to assume that a knowledge graph or a
particular ontology is already the correct representation. Instead, the
literature is being used to determine the representation and reasoning
requirements that the case study must satisfy.

In particular, the following questions remain open:

\begin{enumerate}
    \item Which HAZOP concepts need to be explicitly represented?

    \item Which relation types are required to distinguish causality from
    simple association?

    \item How should multi-step and branching causal chains be represented?

    \item How should multiple causes leading to a common effect be
    represented?

    \item How should safeguards be attached to causal chains?

    \item How should uncertainty in extracted relationships be represented?

    \item Can an extracted causal graph support reasoning over more than one
    causal hop?

    \item Can missing or incomplete causal relationships be identified?
\end{enumerate}

These questions will be used to define the case study and experimental
methodology in the subsequent sections of this week's report.

% ============================================================
% Week 03 --- Part II
% Case-study design and data requirements
% ============================================================

\section{Case-Study Design}

The literature review established that the central objective of this
research is not merely to identify HAZOP entities, but to recover the
causal structure connecting deviations, causes, consequences, and
safeguards. Therefore, the choice of case study is an important part of
the experimental design.

Rather than selecting a process solely because a HAZOP dataset is
available, the case study should provide sufficient causal structure
for evaluating information extraction and subsequent reasoning. The
selected process should also contain enough variation in terminology
and causal descriptions to test whether the proposed NLP/LLM approach
can generalize beyond a small collection of nearly identical sentences.

\subsection{Requirements for the Case Study}

The initial requirements for the case study are summarized in
Figure~\ref{fig:case-study-requirements}.

\begin{figure}[htbp]
\centering
\begin{tikzpicture}[
    node distance=0.45cm,
    box/.style={
        rectangle,
        rounded corners,
        draw=hzblue,
        fill=hzbluebg,
        align=center,
        text width=4.0cm,
        minimum height=0.75cm,
        font=\scriptsize
    },
    arrow/.style={ar}
]

\node[box] (case) {Selected Chemical Process};

\node[box, below left=0.55cm and 0.15cm of case] (reports)
    {Historical HAZOP Reports};

\node[box, below=0.55cm of case] (nodes)
    {Multiple Nodes / Equipment};

\node[box, below right=0.55cm and 0.15cm of case] (chains)
    {Genuine Causal Chains};

\node[box, below=0.55cm of reports] (branch)
    {Branching Causes and Consequences};

\node[box, below=0.55cm of nodes] (safeguards)
    {Safeguards and Recommendations};

\node[box, below=0.55cm of chains] (variation)
    {Linguistic Variation};

\node[box, below=0.55cm of safeguards] (validation)
    {Expert-Verifiable Ground Truth};

\draw[arrow] (case) -- (reports);
\draw[arrow] (case) -- (nodes);
\draw[arrow] (case) -- (chains);

\draw[arrow] (reports) -- (branch);
\draw[arrow] (nodes) -- (safeguards);
\draw[arrow] (chains) -- (variation);

\draw[arrow] (branch) -- (validation);
\draw[arrow] (safeguards) -- (validation);
\draw[arrow] (variation) -- (validation);

\end{tikzpicture}
\caption{Initial requirements for selecting the case study.}
\label{fig:case-study-requirements}
\end{figure}

The following characteristics are considered desirable.

\begin{enumerate}[label=\textbf{\arabic*.}]

    \item \textbf{Availability of HAZOP reports:}
    The case study should have access to sufficiently detailed HAZOP
    worksheets or reports containing deviations, causes, consequences,
    safeguards, and relevant process context.

    \item \textbf{Multiple process nodes:}
    The process should contain multiple equipment items or HAZOP nodes
    so that the proposed method can be evaluated across different
    process contexts.

    \item \textbf{Explicit causal relationships:}
    The reports should contain statements from which causal
    relationships can be recovered rather than only isolated labels
    or keywords.

    \item \textbf{Multi-step causal chains:}
    The data should contain cases in which a cause produces an
    intermediate state or deviation that subsequently produces one or
    more consequences.

    \item \textbf{Branching structure:}
    Ideally, individual deviations should have multiple possible causes
    and consequences. This is important because a simple one-to-one
    mapping would not adequately test causal graph construction.

    \item \textbf{Safeguard information:}
    The HAZOP records should contain safeguards or recommendations so
    that the causal structure can eventually be connected to risk
    mitigation information.

    \item \textbf{Linguistic variation:}
    Different HAZOP entries should express similar causal relationships
    using different terminology or sentence structures. This allows
    the extraction system to be tested for semantic rather than
    purely lexical matching.

    \item \textbf{Expert-verifiable ground truth:}
    The extracted relationships and reconstructed causal chains should
    be capable of being checked against the original HAZOP analysis
    and, where possible, by a process-safety expert.

\end{enumerate}

\subsection{Why These Requirements Matter}

The requirements are motivated by the distinction between
information extraction and causal reasoning identified in the
literature.

For example, consider a simplified process-safety description:

\begin{center}
\begin{tikzpicture}[
    node distance=0.55cm,
    box/.style={
        rectangle,
        rounded corners,
        draw=hzblue,
        fill=hzbluebg,
        align=center,
        text width=3.1cm,
        minimum height=0.65cm,
        font=\scriptsize
    },
    arrow/.style={ar}
]

\node[box] (cause) {Cooling failure};
\node[box, right=0.7cm of cause] (temp) {Temperature increase};
\node[box, right=0.7cm of temp] (pressure) {Pressure increase};
\node[box, right=0.7cm of pressure] (loss) {Loss of containment};

\draw[arrow] (cause) -- (temp);
\draw[arrow] (temp) -- (pressure);
\draw[arrow] (pressure) -- (loss);

\end{tikzpicture}
\end{center}

In this example, identifying the four concepts is only the first
step. A useful system must also recover the relationships between
them. The resulting structure can be represented as

\begin{equation}
C_1 \rightarrow D_1 \rightarrow E_1 \rightarrow C_2,
\end{equation}

where the symbols represent successive causal states rather than
independent HAZOP entities.

Consequently, a case study containing only isolated statements such as
``high pressure'', ``valve failure'', or ``fire'' would be insufficient
for evaluating the main research objective. The dataset must contain
enough contextual information to determine how such concepts are
related.

\subsection{HAZOP Entry as the Basic Unit of Analysis}

The initial experimental unit will be an individual HAZOP entry or
worksheet row. A conventional entry can be viewed as containing
several semantically distinct components:

\begin{center}
\begin{tikzpicture}[
    node distance=0.45cm,
    box/.style={
        rectangle,
        rounded corners,
        draw=hzblue,
        fill=hzbluebg,
        align=center,
        text width=2.4cm,
        minimum height=0.65cm,
        font=\scriptsize
    },
    arrow/.style={ar}
]

\node[box] (node) {Node};
\node[box, right=0.35cm of node] (dev) {Deviation};
\node[box, right=0.35cm of dev] (cause) {Cause};
\node[box, right=0.35cm of cause] (effect) {Consequence};
\node[box, right=0.35cm of effect] (safe) {Safeguard};

\draw[arrow] (node) -- (dev);
\draw[arrow] (dev) -- (cause);
\draw[arrow] (cause) -- (effect);
\draw[arrow] (effect) -- (safe);

\end{tikzpicture}
\end{center}

This representation should not be interpreted as claiming that every
HAZOP entry forms a simple linear chain. Instead, it provides an
initial decomposition of the information contained in a worksheet.
The actual representation may contain multiple causes, intermediate
states, several consequences, and safeguards attached to different
parts of the causal structure.

\subsection{From HAZOP Rows to Causal Propositions}

A major objective of the case-study analysis will therefore be to
transform the textual information in a HAZOP entry into explicit
causal propositions.

For example, a textual entry could conceptually produce propositions
of the form

\begin{align}
\text{Cause} &\xrightarrow{\text{causes}} \text{Deviation},\\
\text{Deviation} &\xrightarrow{\text{leads\_to}} \text{Consequence},\\
\text{Consequence} &\xrightarrow{\text{mitigated\_by}} \text{Safeguard}.
\end{align}

The precise relation vocabulary is intentionally left open at this
stage. Determining whether these relations are sufficient, or whether
additional relations are required to represent intermediate process
states, is one of the questions to be investigated during the
subsequent design phase.

\subsection{Ground-Truth Construction}

The original HAZOP reports will serve as the primary source of
ground-truth information. However, the original worksheet structure
may not directly provide all of the intermediate causal links required
for evaluating multi-hop reasoning.

Therefore, the ground truth will be developed in stages:

\begin{enumerate}[label=\textbf{\arabic*.}]
    \item identify the relevant HAZOP entities in the source entry;
    \item identify explicitly stated causal relationships;
    \item identify multi-step relationships that are supported by the
    text and process context;
    \item represent alternative causes and consequences where they
    exist;
    \item associate safeguards with the corresponding causal scenario;
    \item obtain expert review for relationships that require process
    interpretation.
\end{enumerate}

This distinction is important because an automatically extracted graph
should not be treated as ground truth merely because the model
generated it. The reference structure must originate from the
available HAZOP evidence and, where necessary, expert assessment.

\subsection{Case-Study Selection Strategy}

At this stage, a specific process will not be selected solely on the
basis of convenience. Candidate datasets will instead be evaluated
against the requirements defined above.

The selection can be formulated as a screening process:

\begin{center}
\begin{tikzpicture}[
    node distance=0.55cm,
    box/.style={
        rectangle,
        rounded corners,
        draw=hzblue,
        fill=hzbluebg,
        align=center,
        text width=3.0cm,
        minimum height=0.7cm,
        font=\scriptsize
    },
    arrow/.style={ar}
]

\node[box] (candidate) {Candidate Process / Dataset};

\node[box, below=0.55cm of candidate] (availability)
    {HAZOP Data Available?};

\node[box, below=0.55cm of availability] (structure)
    {Sufficient Causal Structure?};

\node[box, below=0.55cm of structure] (variation)
    {Sufficient Linguistic Variation?};

\node[box, below=0.55cm of variation] (validation)
    {Ground Truth Verifiable?};

\node[box, below=0.55cm of validation] (select)
    {Selected Case Study};

\draw[arrow] (candidate) -- (availability);
\draw[arrow] (availability) -- (structure);
\draw[arrow] (structure) -- (variation);
\draw[arrow] (variation) -- (validation);
\draw[arrow] (validation) -- (select);

\end{tikzpicture}
\end{center}

\subsection{Current Status}

The case-study process is therefore considered an open design
decision at this point. The immediate objective is to identify a
dataset that satisfies the causal and validation requirements rather
than to optimize for dataset size alone.

This approach also allows the research design to remain independent of
a particular process until the available HAZOP material has been
examined. Once candidate reports are obtained, their structure,
terminology, causal depth, branching, and available process context
can be compared against the requirements established in this section.

\subsection{Implications for the Experimental Design}

The case study will eventually support evaluation at several levels:

\begin{center}
\begin{tikzpicture}[
    node distance=0.5cm,
    box/.style={
        rectangle,
        rounded corners,
        draw=hzblue,
        fill=hzbluebg,
        align=center,
        text width=3.2cm,
        minimum height=0.7cm,
        font=\scriptsize
    },
    arrow/.style={ar}
]

\node[box] (text) {HAZOP Text};
\node[box, right=0.7cm of text] (extract)
    {Entity / Relation Extraction};
\node[box, right=0.7cm of extract] (graph)
    {Causal Representation};
\node[box, right=0.7cm of graph] (reason)
    {Multi-hop Reasoning};

\draw[arrow] (text) -- (extract);
\draw[arrow] (extract) -- (graph);
\draw[arrow] (graph) -- (reason);

\end{tikzpicture}
\end{center}

This decomposition will make it possible to distinguish errors caused
by incorrect extraction from errors introduced during graph
construction or causal reasoning. Such separation is necessary for
determining whether an observed failure is an NLP problem, a
representation problem, or a reasoning problem.

The next stage of the research will therefore focus on examining
candidate HAZOP datasets and determining the exact causal schema that
can be reliably annotated from them.

% ============================================================
% Week 03 --- Part III
% Causal Schema, Experimental Design, and Next Steps
% ============================================================

\section{Causal Schema and Experimental Design}

The previous sections established the literature landscape and the
requirements for selecting a suitable HAZOP case study. The remaining
task is to define how the extracted information will be represented,
evaluated, and used for causal reasoning.

\subsection{Initial Causal Representation}

The working representation is based on typed entities and relationships
rather than the original tabular HAZOP format. The initial structure is:

\begin{center}
\begin{tikzpicture}[
    node distance=0.55cm,
    box/.style={
        rectangle,
        rounded corners,
        draw=hzblue,
        fill=hzbluebg,
        align=center,
        text width=2.3cm,
        minimum height=0.65cm,
        font=\scriptsize
    },
    arrow/.style={ar}
]

\node[box] (cause) {Cause};
\node[box, right=0.5cm of cause] (dev) {Deviation};
\node[box, right=0.5cm of dev] (cons) {Consequence};
\node[box, right=0.5cm of cons] (safe) {Safeguard};

\draw[arrow] (cause) -- (dev);
\draw[arrow] (dev) -- (cons);
\draw[arrow] (cons) -- (safe);

\end{tikzpicture}
\end{center}

This should not be interpreted as a restriction to linear chains. A
deviation may have multiple causes and consequences, while a
consequence may itself become an intermediate state in a longer causal
chain.

The representation can therefore be expressed as

\[
G=(V,E),
\]

where \(V\) contains entities and \(E\) contains typed relationships.

The initial entity types are:

\begin{center}
\begin{tabular}{ll}
\textbf{Entity} & \textbf{Description} \\
\hline
Process Node & HAZOP study location \\
Equipment & Physical equipment involved \\
Parameter & Process variable \\
Deviation & HAZOP deviation \\
Cause & Initiating or contributing condition \\
Consequence & Resulting outcome \\
Safeguard & Existing protective measure \\
Recommendation & Proposed protective action
\end{tabular}
\end{center}

The schema is provisional and will be revised after examination of the
selected HAZOP data.

\subsection{Causal Relations}

The initial relation vocabulary is:

\begin{center}
\begin{tabular}{ll}
\textbf{Relation} & \textbf{Meaning} \\
\hline
causes & Cause produces a deviation or state \\
leads\_to & One state produces another state or consequence \\
results\_in & Consequence relation \\
mitigated\_by & Protective measure associated with a scenario \\
associated\_with & Contextual association \\
has\_parameter & Equipment/node--parameter relationship \\
has\_deviation & Node--deviation relationship
\end{tabular}
\end{center}

The final relation vocabulary will be determined from the actual
language and causal structures present in the dataset.

\subsection{Annotation Strategy}

The basic annotation unit will be a causal proposition represented as

\[
(h,r,t),
\]

where \(h\) is the source entity, \(r\) is the relation, and \(t\) is
the target entity.

For example,

\[
(\text{Cooling Failure},\text{causes},\text{High Temperature})
\]

followed by

\[
(\text{High Temperature},\text{leads\_to},\text{Pressure Increase}).
\]

Multiple propositions can then be combined into a causal graph.

The annotation process will distinguish between:

\begin{enumerate}[label=\textbf{\arabic*.}]
    \item entities and process context;
    \item explicitly stated causal relationships;
    \item multi-step relationships;
    \item alternative causes and consequences;
    \item safeguards and recommendations.
\end{enumerate}

Implicit relationships that require substantial process interpretation
will be marked separately rather than automatically treated as
ground truth.

\subsection{Proposed Research Pipeline}

The complete working pipeline for the study is:

\begin{center}
\begin{tikzpicture}[
    node distance=0.4cm,
    box/.style={
        rectangle,
        rounded corners,
        draw=hzblue,
        fill=hzbluebg,
        align=center,
        text width=2.25cm,
        minimum height=0.7cm,
        font=\scriptsize
    },
    arrow/.style={ar}
]

\node[box] (data) {HAZOP Reports};
\node[box, right=0.35cm of data] (extract) {NLP / LLM};
\node[box, right=0.35cm of extract] (causal) {Causal Extraction};
\node[box, right=0.35cm of causal] (graph) {Structured Graph};
\node[box, right=0.35cm of graph] (reason) {Multi-hop Reasoning};

\draw[arrow] (data) -- (extract);
\draw[arrow] (extract) -- (causal);
\draw[arrow] (causal) -- (graph);
\draw[arrow] (graph) -- (reason);

\end{tikzpicture}
\end{center}

The extracted structure will then be compared against the annotated
reference structure and evaluated for causal correctness.

\subsection{Experimental Comparison}

The initial experiments will compare different levels of extraction
and reasoning rather than evaluating only a single model.

\begin{description}[style=nextline,leftmargin=0pt,labelindent=0pt,font=\bfseries]

    \item[Experiment 1: Classical NLP baseline]

    A conventional pipeline using named entity recognition followed by
    relation extraction will provide a baseline for comparison.

    \item[Experiment 2: LLM structured extraction]

    An LLM will be prompted to extract HAZOP entities and typed causal
    relations into a predefined structured format.

    \item[Experiment 3: Explicit causal extraction]

    The extraction task will explicitly require identification of
    causal direction and intermediate states rather than simply
    identifying related entities.

    \item[Experiment 4: Causal graph reasoning]

    The extracted relations will be assembled into a graph and used
    for multi-hop reasoning.

    \item[Experiment 5: Missing-relation experiment]

    Known relationships can be temporarily removed from the graph and
    the system can be tested on whether it can recover the missing
    relationship or causal path.

\end{description}

The last experiment is particularly relevant to the distinction
between extraction and reasoning. A system that merely reproduces
information from the input is different from one that can use an
existing causal structure to identify a missing relationship.

\subsection{Evaluation Metrics}

Evaluation will be performed at multiple levels.

\begin{center}
\begin{tabular}{ll}
\textbf{Level} & \textbf{Possible Metrics} \\
\hline
Entity extraction & Precision, Recall, F1 \\
Relation extraction & Precision, Recall, F1 \\
Causal direction & Relation accuracy \\
Causal chains & Path Precision, Recall, F1 \\
Multi-hop reasoning & 1-hop / 2-hop / 3-hop accuracy \\
Missing relationships & Recovery accuracy, false positives \\
Safety validity & Expert/process-safety assessment
\end{tabular}
\end{center}

This separation is important because a system may achieve good entity
recognition while still producing incorrect causal relationships.

Similarly, correct relation extraction does not necessarily imply
successful multi-hop reasoning.

\subsection{Validation Strategy}

The final evaluation should therefore include both computational and
domain-based validation.

\begin{center}
\begin{tikzpicture}[
    node distance=0.55cm,
    box/.style={
        rectangle,
        rounded corners,
        draw=hzblue,
        fill=hzbluebg,
        align=center,
        text width=3.0cm,
        minimum height=0.7cm,
        font=\scriptsize
    },
    arrow/.style={ar}
]

\node[box] (prediction) {Predicted Causal Structure};

\node[box, below left=0.65cm and 0.35cm of prediction]
    (metric) {Automatic Metrics};

\node[box, below right=0.65cm and 0.35cm of prediction]
    (expert) {Process / Expert Validation};

\node[box, below=1.4cm of prediction]
    (final) {Validated Causal Reasoning};

\draw[arrow] (prediction) -- (metric);
\draw[arrow] (prediction) -- (expert);
\draw[arrow] (metric) -- (final);
\draw[arrow] (expert) -- (final);

\end{tikzpicture}
\end{center}

Automatic metrics will measure agreement with the annotated reference
structure, while expert validation will be necessary for relationships
whose correctness depends on process knowledge or physical
interpretation.

\subsection{Working Research Question}

Based on the literature examined so far, the working research question
is:

\begin{quote}
\textit{Can natural-language processing and large language models
reliably extract the causal structure embedded in HAZOP reports and
represent it in a form that supports explicit multi-hop causal
reasoning?}
\end{quote}

This remains a working research question rather than a finalized claim
of novelty.

The investigation will focus on three supporting questions:

\begin{enumerate}
    \item Can NLP/LLMs reliably extract causal relationships from
    HAZOP text?

    \item Does an explicit causal representation improve multi-hop
    reasoning compared with direct generation of HAZOP conclusions?

    \item Can causal reasoning identify missing or overlooked
    relationships while maintaining process-safety validity?
\end{enumerate}

\subsection{Expected Research Direction}

The current understanding of the project can therefore be summarized
as:

\[
\boxed{\begin{gathered}
\text{Historical HAZOP Reports}\;\rightarrow\;\text{Causal Extraction}\;\rightarrow\\[-1pt]
\text{Structured Representation}\;\rightarrow\;\text{Multi-hop Reasoning}\;\rightarrow\;\text{Validation}
\end{gathered}}
\]

The specific representation---knowledge graph, ontology, causal graph,
or a hybrid approach---will not be fixed until the case-study data and
annotation requirements have been examined.

This preserves flexibility while keeping the central objective
unchanged: recovering and reasoning over the causal structure embedded
in HAZOP documentation.

\subsection{Next Steps}

The immediate next steps are:

\begin{enumerate}[label=\textbf{\arabic*.}]
    \item identify candidate HAZOP datasets;
    \item examine their structure and availability;
    \item select the case study using the requirements defined earlier;
    \item inspect representative HAZOP entries;
    \item refine the entity and relation schema;
    \item define the annotation format;
    \item construct an initial annotated subset;
    \item establish the first extraction baseline.
\end{enumerate}

These steps will provide the transition from literature and research
design into the implementation and experimental phase of the project.
